\documentclass{article}



\begin{document}
Durante o estudo de progressões, é comum encontrar expressões que envolvem termos elevados a potências.
Considere a seguinte sequência definida por uma relação de recorrência:
A sequência an é definida por:
    an+1 = 2an + 3
com a = 1
podemos calcular os primeiros termos passo a passo. Alinhe corretamente as equações a seguir:
    a1 = 1
    a2 = 2a1 + 3 = 5
    a3 = 2a2 + 3 = 13
    a4 = 2a3 + 3 = 29
Tente usar o ambiente align para deixar as igualdades alinhadas no sinal de igual.
Depois represente a forma geral da sequência:
    an = 2n - 3
Por fim, escreva uma frase explicando o que acontece com o valor de 
an conforme  n aumenta, e insira uma fórmula inline que mostre o termo geral aumentado em uma unidade, como por exemplo an+1.    

\end{document}
